\begin{proof}

\begin{enumerate}[label={$\langle1\rangle$\arabic*.},leftmargin=*]
\item Write $\mathscr E=R^1f_*\OO_Y$. Statement \ref{cdp:st:D.1.12} gives its
local freeness, full-fibre base change, and rank two for this
threefold. A local filling $U_i$ retracts onto its central
fibre $S_i$: the retraction of the blown-up base onto
$B_0\cup R$ lifts through its principal compact torus
bundle, by the homotopy lifting property. Therefore (D.25)
identifies
\[
 H^1(S_i;\Z)=H^1(U_i;\Z)
          =\{\lambda\in\Lambda^*:\lambda(w_i)=0\}.
\tag{D.36}\hypertarget{tag-019N}{}\label{cdp:eq:D.36}
\]
The coefficient map into $H^1(U_i,\OO_{U_i})$
corresponds by coherent Leray on the Stein disc to
holomorphic sections of $\mathscr E$. On the punctured
disc, these are constant vectors in $H^1(A,\OO_A)$,
because (D.24) has constant marked periods. The explicit
coefficient map is $\iota$ in \eqref{cdp:eq:D.35}: the $(0,1)$ parts
of $dx,dy$ on $E$ are $d\bar z/2$ and
$\mathrm i\,d\bar z/2$.

\item Each annihilator in \eqref{cdp:eq:D.36} has real dimension three, so
its complex span in the two-dimensional period space is
all of that space. The constant complex relation among
three chosen integral generators holds over the
punctured disc and therefore holds as an identity of
holomorphic sections over the entire disc. At zero,
their three real vectors remain independent: the map
$H^1(S_i;\R)\to H^1(S_i,\OO_{S_i})$ is real-linearly
injective by \ref{cdp:st:D.1.7}, and \ref{cdp:st:D.1.12} identifies this target
with the fibre of $\mathscr E$. Three real independent
vectors cannot lie in one complex line. Their central
complex span thus still has dimension two. Two of the
constant period vectors forming a basis consequently
extend as a frame of $\mathscr E$ across zero.

\item These local frames agree with the same constant period
frame on the smooth base. They therefore glue to
$\mathscr E\simeq\OO^{\oplus2}$ on $\PP^1$.
The cup-product isomorphism in \ref{cdp:st:D.1.12} gives
$R^2f_*\OO_Y\simeq\det\mathscr E\simeq\OO$.
Connected fibres give $f_*\OO_Y=\OO$. Leray on $\PP^1$
proves all four dimensions in \eqref{cdp:eq:D.33}, and evaluation
of the trivial rank-two bundle, followed by the proved
base-change isomorphism, proves coherent surjectivity.

\item The global integral sections are the intersection of
the three annihilators \eqref{cdp:eq:D.36} inside the constant
smooth-fibre lattice. Their defining equations are
$\lambda_3=0$, $\lambda_2+\lambda_4=0$,
$\lambda_1+\lambda_4=0$. The primitive solution is
$(1,1,0,-1)$, giving \eqref{cdp:eq:D.34}. In particular all three
extra fibre components are necessary in the corrected
count: $r=6$, $s=3$, and $t'=4$ give
$6=3-1+4$.

\item Finally $\Pic(Y)=H^1(Y,\OO_Y)$ by the integral vanishing
already proved. Under its coherent restriction
identification with $H^1(A,\OO_A)$, a line bundle is
torsion on a smooth fibre exactly when some positive
integral multiple of its class belongs to
$\iota(\Lambda^*)$, by the exponential sequence.
This is exactly \eqref{cdp:eq:D.35}. At the singular fibre $S_i$
the same exponential sequence and \eqref{cdp:eq:D.36} replace
$\Lambda^*$ by its sublattice $\operatorname{ann}(w_i)$.
That torsion locus is contained in \eqref{cdp:eq:D.35}, so the union
over all fibres is exactly the displayed set. The
real-linear map $\iota:\R^4\to\C^2$ is an isomorphism;
the image of $\Q^4$ is therefore countable and dense.
Its complement is dense and has full Lebesgue measure,
but has no nonempty open subset. This proves the final
assertion and the exact meaning of the generic locus
in this example.

\end{enumerate}
\end{proof}
